\clearpage

\section {Introduction }
\label{sec:intro}

In this technical note, 
we report the results of a 3-flavour $\nu_{\mu}$-disappearance analysis,
performed by the {\em VALOR} group,        
on the combined Run 1+2+3+4 (6.5700$\times$10$^{20}$ POT) dataset.
This dataset was collected in four run periods:
January - June 2010 (Run 1, 0.323$\times$10$^{20}$ POT),
November 2010 - March 2011 (Run 2, 1.108$\times$10$^{20}$ POT),
March - June 2012 (Run 3, 1.579$\times$10$^{20}$ POT, following the recovery work from the damage of the 2011 Tohoku earthquake) and
%in October 2012 - 11th April 2013 (Run 4a, 3.3833$\times$10$^{20}$ POT). 
October 2012 - 8th May 2013 (Run 4, 3.5600$\times$10$^{20}$ POT).
This analysis is an update of the work presented in Refs.\ \cite{OurRun1DisappAnalysis}, \cite{OurRun12DisappAnalysis}, 
\cite{OurRun12DisappAnalysis3Flavour},
\cite{OurRun123DisappAnalysis3Flavour} and 
\cite{OurRun1234aDisappAnalysis3Flavour} with improved methods and an enlarged dataset. 

During our Run 1+2+3 analysis, it was found that the choice of octant of $\theta_{23}$ makes a significant difference to the 
confidence region in the $\sin^{2}(2\theta_{23})$ - $|\Delta m^{2}_{32}|$ oscillation parameter space. For this reason, best-fit values 
were found and confidence regions constructed using $\sin^{2}\theta_{23}$ as the mixing angle instead of $\sin^{2}(2\theta_{23})$ 
\cite{OurRun123DisappAnalysis3Flavour}.   
This use of $\sin^{2}\theta_{23}$ is continued in this Run 1+2+3+4 analysis. There has been a change in the mass-squared 
differences used in this analysis: we now use $|\Delta m^{2}_{32}|$ for the normal mass hierarchy and $|\Delta m^{2}_{31}|$ for 
the inverted hierarchy, whereas, in our Run 1+2+3 analysis, we used $|\Delta m^{2}_{32}|$ for both hierarchies. This change has 
been made because $|\Delta m^{2}_{32}|$ is the largest mass-squared difference in the inverted hierarchy, but the second largest 
in the normal hierarchy.
% The inputs to the calculation of the oscillation probabilities are +$|\Delta m^{2}_{32}|$ for the normal 
% hierarchy and -$|\Delta m^{2}_{31}|$ for the inverted hierarchy.

These oscillation parameters are determined
by fitting the predicted reconstructed energy spectrum
of 1 $\mu$-like ring events at SuperK to the observed spectrum.
The method of calculation of this predicted spectrum is presented in Sec.\ \ref{sec:spectrum_prediction}.
This analysis uses flux tuning based on recent NA61 data \cite{FluxTuning}\cite{NA61T2K}. 
It also uses an indirect extrapolation from ND280 to SuperK, which is implemented via a fit of flux and cross-section parameters to
ND280 measurements of the muon momentum-angle distributions of $\nu_\mu$
CC-0-pion-like, CC-1-pion-like and CC-other events whose 
vertices are in FGD1 \cite{BANFF2013}.
Updated nominal predictions of the reconstructed energy spectrum of single $\mu$-like ring events in SuperK are presented for Run 1+2+3+4
in Sec.\ \ref{sec:spectrum_prediction}; the changes to the analysis summarized in Sec.\ \ref{sec:changes} are incorporated 
into these predictions. 

We consider a total of 41 energy scale, neutrino flux, neutrino cross-section,
efficiency, initial-state nuclear environment and final-state rescattering systematics (nuisance parameters).
The oscillation parameters $\sin^{2}\theta_{12}$, $\sin^{2}\theta_{13}$, $\Delta m^{2}_{21}$ and $\delta_{CP}$ are also considered 
as nuisance parameters in this analysis. These nuisance parameters and their effects 
on the extrapolated single $\mu$-like ring event reconstructed energy spectrum are described in Sec.\ \ref{sec:systematics}.
Further details of them are also given in a separate technical note
supplement \cite{OurRun1234DisappAnalysisSupplement}.

The oscillation analysis is performed using neutrino oscillation probabilities computed in a 3-flavour
framework including constant-density matter effects. 
Details of our oscillation probability calculations, estimates of the numerical accuracy of our
calculations and comparisons with calculations used by the alternative T2K oscillation analyses 
are given in Appendices A and B of Ref. \cite{OurRun123DisappAnalysis3Flavour}.

As in our previous analyses, the oscillation fit is made using the likelihood-ratio method \cite{StatAnalysisCowan}.
The mathematical details of this method are given in Sec.\ 3.2 in Ref.\ \cite{OurRun1DisappAnalysis}.
In this Run 1+2+3+4 analysis, we allow all 45 nuisance parameters (41 systematic parameters 
and four non-23-sector oscillation parameters) to float.
The oscillation fitting method is presented in detail in Sec.\ \ref{sec:oscfit_method},
while the chosen goodness-of-fit test, suitable for the current statistics, is described in Sec.\ \ref{sec:gof_test}.
The performance of the fitter was tested in detail for the Run 1+2+3+4 fit using ensembles of toy MC experiments, and
a summary of the results of these tests is presented in Sec.\ \ref{sec:validation}.

The 2D confidence regions in $\sin^{2}\theta_{23}$ and $|\Delta m^{2}_{32}|$ or $|\Delta m^{2}_{31}|$ are constructed using 
the Feldman-Cousins method \cite{FeldmanCousins} as in our previous analyses. 
Systematic errors are included in the oscillation contours using the 
Cousins-Highland method (Bayesian integration) \cite{CousinsHighland}.
Confidence contours constructed with the constant-$\Delta\chi^{2}$ method are also shown for both fits for comparison.
Both of these methods are summarized in Sec.\ \ref{sec:contour_construction}, and
the exact methodology used for incorporating systematic errors is explained in detail.
1D confidence intervals are also calculated for each of $\sin^{2}\theta_{23}$ and $|\Delta m^{2}_{32}|$ or $|\Delta m^{2}_{31}|$ separately 
using a new procedure that is an extension of the Feldman and Cousins method; this procedure is described in Sec.\ 
\ref{sec:contour_construction}.

This analysis is similar to our Run 1+2+3 analysis with higher statistics, though certain analysis inputs and fitting methods 
have been updated. For the benefit of readers, all changes with respect to our previously published 
Run 1+2+3 $\nu_{\mu}$-disappearance analysis described in Ref.\ \cite{OurRun123DisappAnalysis3Flavour} 
are summarized in Sec.\ \ref{sec:changes}.

The results of the measurements of $\sin^{2}\theta_{23}$ and $|\Delta m^{2}_{32}|$ or $|\Delta m^{2}_{31}|$ 
using the Run 1+2+3+4 dataset are given in Sec.\ \ref{sec:results_bestfit} of this note, 
while the goodness-of-fit tests are presented in Sec.\ \ref{sec:results_gof}.
The best-fit spectra are shown in Sec.\ \ref{sec:results_bestfit_spectra} for each of the normal and inverted mass hierarchy assumptions for the 
Run 1+2+3+4 dataset.
The Feldman-Cousins and constant-$\Delta\chi^{2}$ confidence regions are shown in Sec.\ \ref{sec:results_contours}. 
Comparisons of our results obtained between different methods for constructing confidence regions are presented in Sec.\ \ref{sec:results_contour_comparisons}.
In that section, we also show comparisons of our T2K Run 1+2+3+4 results with results from MINOS and SuperK
as well as with the expected T2K sensitivity for an integrated exposure of 6.5700$\times$10$^{20}$ POT.

Finally, a summary is given in Sec.\ \ref{sec:summary}.  

\clearpage

